\begin{filecontents*}[overwrite]{custom.bib}
% Use this file for citations not found in the ACL Anthology (contained in "anthology.bib").

@book{Aho:72,
    author  = {Alfred V. Aho and Jeffrey D. Ullman},
    title   = {The Theory of Parsing, Translation and Compiling},
    year    = "1972",
    volume  = "1",
    publisher = {Prentice-Hall},
    address = {Englewood Cliffs, NJ}
}

@book{APA:83,
    author  = {{American Psychological Association}},
    title   = {Publications Manual},
    year    = "1983",
    publisher = {American Psychological Association},
    address = {Washington, DC}
}

@article{Chandra:81,
	author = {Ashok K. Chandra and Dexter C. Kozen and Larry J. Stockmeyer},
	year = "1981",
	title = {Alternation},
	journal = {Journal of the Association for Computing Machinery},
	volume = "28",
	number = "1",
	pages = "114--133",
	doi = "10.1145/322234.322243",
}

@inproceedings{andrew2007scalable,
  title={Scalable training of {L1}-regularized log-linear models},
  author={Andrew, Galen and Gao, Jianfeng},
  booktitle={Proceedings of the 24th International Conference on Machine Learning},
  pages={33--40},
  year={2007},
}

@book{Gusfield:97,
    author  = {Dan Gusfield},
    title   = {Algorithms on Strings, Trees and Sequences},
    year    = "1997",
    publisher = {Cambridge University Press},
    address = {Cambridge, UK}
}

@article{rasooli-tetrault-2015,
    author    = {Mohammad Sadegh Rasooli and Joel R. Tetreault},
    title     = {Yara Parser: {A} Fast and Accurate Dependency Parser},
    journal   = {Computing Research Repository},
    volume    = {arXiv:1503.06733},
    year      = {2015},
    url       = {http://arxiv.org/abs/1503.06733},
    note    = {version 2}
}

@article{Ando2005,
	Acmid = {1194905},
	Author = {Ando, Rie Kubota and Zhang, Tong},
	Issn = {1532-4435},
	Issue_Date = {12/1/2005},
	Journal = {Journal of Machine Learning Research},
	Month = dec,
	Numpages = {37},
	Pages = {1817--1853},
	Publisher = {JMLR.org},
	Title = {A Framework for Learning Predictive Structures from Multiple Tasks and Unlabeled Data},
	Volume = {6},
	Year = {2005}
}

@inproceedings{shardlow-etal-2021-semeval,
    title = "{S}em{E}val-2021 Task 1: Lexical Complexity Prediction",
    author = "Shardlow, Matthew  and
      Evans, Richard  and
      Paetzold, Gustavo Henrique  and
      Zampieri, Marcos",
    editor = "Palmer, Alexis  and
      Schneider, Nathan  and
      Schluter, Natalie  and
      Emerson, Guy  and
      Herbelot, Aurelie  and
      Zhu, Xiaodan",
    booktitle = "Proceedings of the 15th International Workshop on Semantic Evaluation (SemEval-2021)",
    month = aug,
    year = "2021",
    address = "Online",
    publisher = "Association for Computational Linguistics",
    url = "https://aclanthology.org/2021.semeval-1.1/",
    doi = "10.18653/v1/2021.semeval-1.1",
    pages = "1--16"
}

@inproceedings{yimam-etal-2018-report,
    title = "A Report on the Complex Word Identification Shared Task 2018",
    author = {Yimam, Seid Muhie  and
      Biemann, Chris  and
      Malmasi, Shervin  and
      Paetzold, Gustavo  and
      Specia, Lucia  and
      {\v{S}}tajner, Sanja  and
      Tack, Ana{\"i}s  and
      Zampieri, Marcos},
    editor = "Tetreault, Joel  and
      Burstein, Jill  and
      Kochmar, Ekaterina  and
      Leacock, Claudia  and
      Yannakoudakis, Helen",
    booktitle = "Proceedings of the Thirteenth Workshop on Innovative Use of {NLP} for Building Educational Applications",
    month = jun,
    year = "2018",
    address = "New Orleans, Louisiana",
    publisher = "Association for Computational Linguistics",
    url = "https://aclanthology.org/W18-0507/",
    doi = "10.18653/v1/W18-0507",
    pages = "66--78"
}

@inproceedings{shardlow-etal-2024-multilingual,
    title = "Multilingual Resources for Lexical Complexity Prediction: A Review",
    author = "Shardlow, Matthew  and
      North, Kai  and
      Zampieri, Marcos",
    editor = "Nunzio, Giorgio Maria Di  and
      Vezzani, Federica  and
      Ermakova, Liana  and
      Azarbonyad, Hosein  and
      Kamps, Jaap",
    booktitle = "Proceedings of the Workshop on DeTermIt! Evaluating Text Difficulty in a Multilingual Context @ LREC-COLING 2024",
    month = may,
    year = "2024",
    address = "Torino, Italia",
    publisher = "ELRA and ICCL",
    url = "https://aclanthology.org/2024.determit-1.5/",
    pages = "51--59"
}

@inproceedings{shardlow-etal-2024-bea,
    title = "The {BEA} 2024 Shared Task on the Multilingual Lexical Simplification Pipeline",
    author = {Shardlow, Matthew  and
      Alva-Manchego, Fernando  and
      Batista-Navarro, Riza  and
      Bott, Stefan  and
      Calderon Ramirez, Saul  and
      Cardon, R{\'e}mi  and
      Fran{\c{c}}ois, Thomas  and
      Hayakawa, Akio  and
      Horbach, Andrea  and
      H{\"u}lsing, Anna  and
      Ide, Yusuke  and
      Imperial, Joseph Marvin  and
      Nohejl, Adam  and
      North, Kai  and
      Occhipinti, Laura  and
      Rojas, Nelson Per{\'e}z  and
      Raihan, Nishat  and
      Ranasinghe, Tharindu  and
      Salazar, Martin Solis  and
      {\v{S}}tajner, Sanja  and
      Zampieri, Marcos  and
      Saggion, Horacio},
    editor = {Kochmar, Ekaterina  and
      Bexte, Marie  and
      Burstein, Jill  and
      Horbach, Andrea  and
      Laarmann-Quante, Ronja  and
      Tack, Ana{\"i}s  and
      Yaneva, Victoria  and
      Yuan, Zheng},
    booktitle = "Proceedings of the 19th Workshop on Innovative Use of NLP for Building Educational Applications (BEA 2024)",
    month = jun,
    year = "2024",
    address = "Mexico City, Mexico",
    publisher = "Association for Computational Linguistics",
    url = "https://aclanthology.org/2024.bea-1.51/",
    pages = "571--589"
}

@inproceedings{qiu-etal-2024-complex,
    title = "{C}omp{L}ex-{ZH}: A New Dataset for Lexical Complexity Prediction in {M}andarin and {C}antonese",
    author = "Qiu, Le  and
      Guo, Shanyue  and
      Wong, Tak-Sum  and
      Chersoni, Emmanuele  and
      Lee, John  and
      Huang, Chu-Ren",
    editor = "Shardlow, Matthew  and
      Saggion, Horacio  and
      Alva-Manchego, Fernando  and
      Zampieri, Marcos  and
      North, Kai  and
      {\v{S}}tajner, Sanja  and
      Stodden, Regina",
    booktitle = "Proceedings of the Third Workshop on Text Simplification, Accessibility and Readability (TSAR 2024)",
    month = nov,
    year = "2024",
    address = "Miami, Florida, USA",
    publisher = "Association for Computational Linguistics",
    url = "https://aclanthology.org/2024.tsar-1.3/",
    doi = "10.18653/v1/2024.tsar-1.3",
    pages = "20--26"
}

@inproceedings{saggion-etal-2024-lexical,
    title = "Lexical Complexity Prediction and Lexical Simplification for {C}atalan and {S}panish: Resource Creation, Quality Assessment, and Ethical Considerations",
    author = "Saggion, Horacio  and
      Bott, Stefan  and
      Szasz, Sandra  and
      P{\'e}rez, Nelson  and
      Calder{\'o}n, Sa{\'u}l  and
      Sol{\'i}s, Mart{\'i}n",
    editor = "Shardlow, Matthew  and
      Saggion, Horacio  and
      Alva-Manchego, Fernando  and
      Zampieri, Marcos  and
      North, Kai  and
      {\v{S}}tajner, Sanja  and
      Stodden, Regina",
    booktitle = "Proceedings of the Third Workshop on Text Simplification, Accessibility and Readability (TSAR 2024)",
    month = nov,
    year = "2024",
    address = "Miami, Florida, USA",
    publisher = "Association for Computational Linguistics",
    url = "https://aclanthology.org/2024.tsar-1.9/",
    doi = "10.18653/v1/2024.tsar-1.9",
    pages = "82--94"
}

@inproceedings{conneau2020xlmr,
  title        = {Unsupervised Cross-lingual Representation Learning at Scale},
  author       = {Conneau, Alexis and Khandelwal, Kartikay and Goyal, Naman and Chaudhary, Vishrav and Wenzek, Guillaume and Guzm{\'a}n, Francisco and Grave, Edouard and Ott, Myle and Zettlemoyer, Luke and Stoyanov, Veselin},
  booktitle    = {Proceedings of the 58th Annual Meeting of the Association for Computational Linguistics (ACL)},
  pages        = {8440--8451},
  year         = {2020},
  url          = {https://aclanthology.org/2020.acl-main.747}
}

@inproceedings{pires-etal-2019-multilingual,
    title = "How Multilingual is Multilingual {BERT}?",
    author = "Pires, Telmo  and
      Schlinger, Eva  and
      Garrette, Dan",
    editor = "Korhonen, Anna  and
      Traum, David  and
      M{\`a}rquez, Llu{\'i}s",
    booktitle = "Proceedings of the 57th Annual Meeting of the Association for Computational Linguistics",
    month = jul,
    year = "2019",
    address = "Florence, Italy",
    publisher = "Association for Computational Linguistics",
    url = "https://aclanthology.org/P19-1493/",
    doi = "10.18653/v1/P19-1493",
    pages = "4996--5001"
}

@inproceedings{wu-dredze-2020-languages,
    title = "Are All Languages Created Equal in Multilingual {BERT}?",
    author = "Wu, Shijie  and
      Dredze, Mark",
    editor = "Gella, Spandana  and
      Welbl, Johannes  and
      Rei, Marek  and
      Petroni, Fabio  and
      Lewis, Patrick  and
      Strubell, Emma  and
      Seo, Minjoon  and
      Hajishirzi, Hannaneh",
    booktitle = "Proceedings of the 5th Workshop on Representation Learning for NLP",
    month = jul,
    year = "2020",
    address = "Online",
    publisher = "Association for Computational Linguistics",
    url = "https://aclanthology.org/2020.repl4nlp-1.16/",
    doi = "10.18653/v1/2020.repl4nlp-1.16",
    pages = "120--130"
}

@inproceedings{peters-etal-2019-tune,
    title = "To Tune or Not to Tune? Adapting Pretrained Representations to Diverse Tasks",
    author = "Peters, Matthew E.  and
      Ruder, Sebastian  and
      Smith, Noah A.",
    editor = "Augenstein, Isabelle  and
      Gella, Spandana  and
      Ruder, Sebastian  and
      Kann, Katharina  and
      Can, Burcu  and
      Welbl, Johannes  and
      Conneau, Alexis  and
      Ren, Xiang  and
      Rei, Marek",
    booktitle = "Proceedings of the 4th Workshop on Representation Learning for NLP (RepL4NLP-2019)",
    month = aug,
    year = "2019",
    address = "Florence, Italy",
    publisher = "Association for Computational Linguistics",
    url = "https://aclanthology.org/W19-4302/",
    doi = "10.18653/v1/W19-4302",
    pages = "7--14"
}

@inproceedings{hale2001probabilistic,
  title        = {A Probabilistic {Earley} Parser as a Psycholinguistic Model},
  author       = {Hale, John},
  booktitle    = {Proceedings of the Second Meeting of the North American Chapter of the Association for Computational Linguistics (NAACL)},
  year         = {2001},
  url          = {https://aclanthology.org/N01-1021}
}

@article{rayner1998eye,
  title        = {Eye movements in reading and information processing: 20 years of research},
  author       = {Rayner, Keith},
  journal      = {Psychological Bulletin},
  volume       = {124},
  number       = {3},
  pages        = {372--422},
  year         = {1998},
  doi          = {10.1037/0033-2909.124.3.372}
}

@inproceedings{hollenstein-etal-2021-multilingual,
    title = "Multilingual Language Models Predict Human Reading Behavior",
    author = {Hollenstein, Nora  and
      Pirovano, Federico  and
      Zhang, Ce  and
      J{\"a}ger, Lena  and
      Beinborn, Lisa},
    editor = "Toutanova, Kristina  and
      Rumshisky, Anna  and
      Zettlemoyer, Luke  and
      Hakkani-Tur, Dilek  and
      Beltagy, Iz  and
      Bethard, Steven  and
      Cotterell, Ryan  and
      Chakraborty, Tanmoy  and
      Zhou, Yichao",
    booktitle = "Proceedings of the 2021 Conference of the North American Chapter of the Association for Computational Linguistics: Human Language Technologies",
    month = jun,
    year = "2021",
    address = "Online",
    publisher = "Association for Computational Linguistics",
    url = "https://aclanthology.org/2021.naacl-main.10/",
    doi = "10.18653/v1/2021.naacl-main.10",
    pages = "106--123"
}

@inproceedings{alves-2025-benchmarking,
    title = "Benchmarking Language Model Surprisal for Eye-Tracking Predictions in {B}razilian {P}ortuguese",
    author = "Alves, Diego",
    editor = "Acarturk, Cengiz  and
      Nasir, Jamal  and
      Can, Burcu  and
      Coltekin, Cagr{\i}",
    booktitle = "Proceedings of the First International Workshop on Gaze Data and Natural Language Processing",
    month = sep,
    year = "2025",
    address = "Varna, Bulgaria",
    publisher = "INCOMA Ltd., Shoumen, BULGARIA",
    url = "https://aclanthology.org/2025.gaze4nlp-1.2/",
    pages = "7--17"
}

@inproceedings{gorman-bedrick-2019-need,
    title = "We Need to Talk about Standard Splits",
    author = "Gorman, Kyle  and
      Bedrick, Steven",
    editor = "Korhonen, Anna  and
      Traum, David  and
      M{\`a}rquez, Llu{\'i}s",
    booktitle = "Proceedings of the 57th Annual Meeting of the Association for Computational Linguistics",
    month = jul,
    year = "2019",
    address = "Florence, Italy",
    publisher = "Association for Computational Linguistics",
    url = "https://aclanthology.org/P19-1267/",
    doi = "10.18653/v1/P19-1267",
    pages = "2786--2791"
}

@inproceedings{lundberg2017shap,
  title        = {A Unified Approach to Interpreting Model Predictions},
  author       = {Lundberg, Scott M. and Lee, Su-In},
  booktitle    = {Advances in Neural Information Processing Systems (NeurIPS)},
  pages        = {4765--4774},
  year         = {2017},
  url          = {https://arxiv.org/abs/1705.07874}
}

@inproceedings{wei2022cot,
  title        = {Chain-of-Thought Prompting Elicits Reasoning in Large Language Models},
  author       = {Wei, Jason and Wang, Xuezhi and Schuurmans, Dale and Bosma, Maarten and Ichter, Brian and Xia, Fei and Chi, Ed H. and Le, Quoc V. and Zhou, Denny},
  booktitle    = {Advances in Neural Information Processing Systems (NeurIPS)},
  year         = {2022},
  url          = {https://arxiv.org/abs/2201.11903}
}

@inproceedings{masala-etal-2024-vorbesti,
    title = "``Vorbe{\c s}ti Rom{\^a}ne{\c s}te?'' A Recipe to Train Powerful {R}omanian {LLM}s with {E}nglish Instructions",
    author = "Masala, Mihai  and
      Ilie-Ablachim, Denis  and
      Dima, Alexandru  and
      Corlatescu, Dragos Georgian  and
      Zavelca, Miruna-Andreea  and
      Olaru, Ovio  and
      Terian, Simina-Maria  and
      Terian, Andrei  and
      Leordeanu, Marius  and
      Velicu, Horia  and
      Popescu, Marius  and
      Dascalu, Mihai  and
      Rebedea, Traian",
    editor = "Al-Onaizan, Yaser  and
      Bansal, Mohit  and
      Chen, Yun-Nung",
    booktitle = "Findings of the Association for Computational Linguistics: EMNLP 2024",
    month = nov,
    year = "2024",
    address = "Miami, Florida, USA",
    publisher = "Association for Computational Linguistics",
    url = "https://aclanthology.org/2024.findings-emnlp.681/",
    doi = "10.18653/v1/2024.findings-emnlp.681",
    pages = "11632--11647"
}
\end{filecontents*}

\documentclass[11pt]{article}

% Change "review" to "final" to generate the final (sometimes called camera-ready) version.
% Change to "preprint" to generate a non-anonymous version with page numbers.
\usepackage[review]{acl}

% Standard package includes
\usepackage{times}
\usepackage{latexsym}

% For proper rendering and hyphenation of words containing Latin characters (including in bib files)
\usepackage[T1]{fontenc}
% For Vietnamese characters
% \usepackage[T5]{fontenc}
% See https://www.latex-project.org/help/documentation/encguide.pdf for other character sets

% This assumes your files are encoded as UTF8
\usepackage[utf8]{inputenc}

% This is not strictly necessary, and may be commented out,
% but it will improve the layout of the manuscript,
% and will typically save some space.
\usepackage{microtype}

% This is also not strictly necessary, and may be commented out.
% However, it will improve the aesthetics of text in
% the typewriter font.
\usepackage{inconsolata}

%Including images in your LaTeX document requires adding
%additional package(s)
\usepackage{graphicx}

% If the title and author information does not fit in the area allocated, uncomment the following
%
%\setlength\titlebox{<dim>}
%
% and set <dim> to something 5cm or larger.

\title{Instructions for *ACL Proceedings}

% Author information can be set in various styles:
% For several authors from the same institution:
% \author{Author 1 \and ... \and Author n \\
%         Address line \\ ... \\ Address line}
% if the names do not fit well on one line use
%         Author 1 \\ {\bf Author 2} \\ ... \\ {\bf Author n} \\
% For authors from different institutions:
% \author{Author 1 \\ Address line \\  ... \\ Address line
%         \And  ... \And
%         Author n \\ Address line \\ ... \\ Address line}
% To start a separate ``row'' of authors use \AND, as in
% \author{Author 1 \\ Address line \\  ... \\ Address line
%         \AND
%         Author 2 \\ Address line \\ ... \\ Address line \And
%         Author 3 \\ Address line \\ ... \\ Address line}

\author{First Author \\
  Affiliation / Address line 1 \\
  Affiliation / Address line 2 \\
  Affiliation / Address line 3 \\
  \texttt{email@domain} \\\And
  Second Author \\
  Affiliation / Address line 1 \\
  Affiliation / Address line 2 \\
  Affiliation / Address line 3 \\
  \texttt{email@domain} \\}

%\author{
%  \textbf{First Author\textsuperscript{1}},
%  \textbf{Second Author\textsuperscript{1,2}},
%  \textbf{Third T. Author\textsuperscript{1}},
%  \textbf{Fourth Author\textsuperscript{1}},
%\\
%  \textbf{Fifth Author\textsuperscript{1,2}},
%  \textbf{Sixth Author\textsuperscript{1}},
%  \textbf{Seventh Author\textsuperscript{1}},
%  \textbf{Eighth Author \textsuperscript{1,2,3,4}},
%\\
%  \textbf{Ninth Author\textsuperscript{1}},
%  \textbf{Tenth Author\textsuperscript{1}},
%  \textbf{Eleventh E. Author\textsuperscript{1,2,3,4,5}},
%  \textbf{Twelfth Author\textsuperscript{1}},
%\\
%  \textbf{Thirteenth Author\textsuperscript{3}},
%  \textbf{Fourteenth F. Author\textsuperscript{2,4}},
%  \textbf{Fifteenth Author\textsuperscript{1}},
%  \textbf{Sixteenth Author\textsuperscript{1}},
%\\
%  \textbf{Seventeenth S. Author\textsuperscript{4,5}},
%  \textbf{Eighteenth Author\textsuperscript{3,4}},
%  \textbf{Nineteenth N. Author\textsuperscript{2,5}},
%  \textbf{Twentieth Author\textsuperscript{1}}
%\\
%\\
%  \textsuperscript{1}Affiliation 1,
%  \textsuperscript{2}Affiliation 2,
%  \textsuperscript{3}Affiliation 3,
%  \textsuperscript{4}Affiliation 4,
%  \textsuperscript{5}Affiliation 5
%\\
%  \small{
%    \textbf{Correspondence:} \href{mailto:email@domain}{email@domain}
%  }
%}

\begin{document}
\maketitle
\begin{abstract}
This document is a supplement to the general instructions for *ACL authors. It contains instructions for using the \LaTeX{} style files for ACL conferences.
The document itself conforms to its own specifications, and is therefore an example of what your manuscript should look like.
These instructions should be used both for papers submitted for review and for final versions of accepted papers.
\end{abstract}

\makeatletter
\def\pinlabel#1#2{\expandafter\gdef\csname r@#1\endcsname{{#2}{1}{}{}{}}}
\pinlabel{section:approach}{3}
\pinlabel{sec:tuning}{3.5}
\pinlabel{sec:results}{4}
\pinlabel{sec:negative}{5}
\makeatother

\section{Introduction}
\label{section:intro}

Making a text usable by a given audience depends, before anything else,
on its vocabulary. A word that is routine for a specialist can stop a
language learner or a reader from outside the field, taking the
surrounding sentence with it. Deciding which
words will cause this, and for whom, is what lexical complexity
prediction (LCP) does: it assigns a word in context a score on a
continuous scale \citep{shardlow-etal-2021-semeval}. Text simplification
uses those scores to choose what to replace, readability assessment to
aggregate difficulty over a document, and language-teaching material to
keep vocabulary within a learner's range; complexity prediction is for
this reason typically the first stage of a simplification pipeline
\citep{qiu-etal-2024-complex}.

Nearly every LCP system is built the same way: a pretrained encoder is
fine-tuned on complexity annotations in the language of interest. The
recipe works wherever those annotations exist, and it quietly assumes the
one resource most languages do not have. Surveying the datasets published
between 2018 and 2024, \citet{shardlow-etal-2024-multilingual} find nine
of them, covering eight languages. A language outside that set has
nothing to fine-tune on, and buying its way in means paying for a new
annotation effort of its own.

Our target language is Romanian, and we use no Romanian complexity
labels at any point, whether for training, validation or model
selection. Working without them means no design decision can be checked
on Romanian itself, so every choice has to be justified on other
languages. Romanian adds a problem of its own: the most reliable
predictor of word difficulty misfires on it. Word frequency governs how
long readers fixate a word \citep{rayner1998eye} and ranks first in our
own attribution analysis (Section~\ref{sec:results}), but Romanian
inflects heavily and the frequency lists available for it cover
inflected forms poorly: of the 1587 Romanian target tokens in our test
set, $11.7\%$ are absent from the list and would be scored as maximally
rare, and for a further $36.0\%$ the listed form is rarer than its own
lemma. A method for Romanian therefore needs complexity knowledge
imported from other languages, and a frequency estimate that survives
inflection.

In this paper, we predict the lexical complexity of Romanian words in
context without a single Romanian training label, and we make four
contributions. First, we combine four label-free signals into a single
predictor: cross-lingual transfer through a frozen multilingual encoder,
a language-model complexity rating, surprisal, and a reading-effort
score derived from eye-tracking. It scores $69.74$ on our test set,
under a metric that combines Pearson correlation with $R^2$.
Second, we introduce a lemma-frequency correction for
morphologically rich languages, substituting the lemma frequency when the
inflected form is missing or rarer; it is worth $+4.09$, the largest
single gain we measure, and applies to any inflected language served
by the same frequency lists. Third, we show that the four signals are
complementary rather than redundant: the cross-lingual signal acts mainly
as a frequency detector, the rating and surprisal are near-independent,
and eye-tracking captures a difficulty the others miss, the effort of
fitting a word into its context, which shows up only in \emph{late} gaze
measures.
Fourth, we report two negative results: fine-tuning the encoder, full or
LoRA, is consistently \emph{worse} than leaving it frozen, and in-source
validation is an unreliable guide, since three configurations that win on
English-to-Romance transfer lose on Romanian.

Our setting differs from the LCP literature in the resource it assumes
rather than in the model it uses: where that literature trains on
target-language annotations, we treat complexity as knowledge to be
imported. We also invert two conventions: gaze becomes evidence about
complexity rather than something to predict from text, and the encoder
stays frozen rather than fine-tuned. Section~\ref{sec:related} sets all
three differences against the work they concern.

The rest of the paper describes the four signals and their fusion
(Section~\ref{section:approach}), reports the ablation and the analysis
(Section~\ref{sec:results}), and collects what did not work
(Section~\ref{sec:negative}).

\section{Related Work}
\label{sec:related}

\paragraph{Lexical complexity prediction.}
Annotated complexity data exists for only a handful of languages, and
Romanian is not among them. The task itself began as classification over
annotated target words \citep{yimam-etal-2018-report}; SemEval-2021
Task~1 replaced it with a continuous scale on the CompLex corpus
\citep{shardlow-etal-2021-semeval}, which is the formulation we adopt.
Zero-shot evaluation is not itself new: the 2018 task already had a
multilingual track whose French test set came with no French training
data, though it asked only for binary classification and supplied
training data in three languages, one of them Romance. Coverage has since
grown to
Catalan, Spanish, Mandarin and Cantonese
\citep{saggion-etal-2024-lexical, qiu-etal-2024-complex} and to the ten
languages of the MLSP shared task \citep{shardlow-etal-2024-bea}, but the
review of \citet{shardlow-etal-2024-multilingual} shows annotated LCP
data still concentrated in a few high-resource languages. Romanian
appears in none of these resources. We use the English and Romance
portions of them as the source side of our transfer.

\paragraph{Cross-lingual transfer without target labels.}
Multilingual encoders such as XLM-R \citep{conneau2020xlmr} place
translation-equivalent words close together, which makes zero-shot
transfer of a regressor trained on one language plausible. Transfer
quality is far from uniform across languages, however, and degrades for
languages with less pretraining data and richer morphology
\citep{pires-etal-2019-multilingual, wu-dredze-2020-languages}, which is
precisely the regime Romanian occupies. Transfer is moreover normally
realised by fine-tuning the encoder on the source language; that feature
extraction can outperform fine-tuning when the target task is distant
from the source is known \citep{peters-etal-2019-tune}, but we are aware
of no prior report of it for cross-lingual complexity transfer. This
motivates our frozen-encoder design, and the lemma-frequency correction
completes it: the encoder supplies cross-lingual semantics, the
correction repairs the lexical feature that morphology corrupts.

\paragraph{Surprisal and eye movements as difficulty signals.}
Surprisal theory predicts that processing effort scales with the
negative log probability of a word in context \citep{hale2001probabilistic},
and gaze measures such as longer fixations, re-reading and regressions
are the standard behavioural correlate of that effort
\citep{rayner1998eye}. Language-model surprisal predicts reading times
across languages \citep{hollenstein-etal-2021-multilingual}, including
in Romance languages \citep{alves-2025-benchmarking}. Both literatures
run from text to behaviour; we take the converse view, reading surprisal
and gaze as label-free evidence about complexity.

\paragraph{Evaluation without a development set.}
Selecting a model against the same data used to report its score inflates
that score \citep{gorman-bedrick-2019-need}. With no development set to
tune on, we state our selection procedure explicitly and quantify its
effect (Section~\ref{sec:tuning}).

\section{Engines}

To produce a PDF file, pdf\LaTeX{} is strongly recommended (over original \LaTeX{} plus dvips+ps2pdf or dvipdf).
The style file \texttt{acl.sty} can also be used with
lua\LaTeX{} and
Xe\LaTeX{}, which are especially suitable for text in non-Latin scripts.
The file \texttt{acl\_lualatex.tex} in this repository provides
an example of how to use \texttt{acl.sty} with either
lua\LaTeX{} or
Xe\LaTeX{}.

\section{Preamble}

The first line of the file must be
\begin{quote}
\begin{verbatim}
\documentclass[11pt]{article}
\end{verbatim}
\end{quote}

To load the style file in the review version:
\begin{quote}
\begin{verbatim}
\usepackage[review]{acl}
\end{verbatim}
\end{quote}
For the final version, omit the \verb|review| option:
\begin{quote}
\begin{verbatim}
\usepackage{acl}
\end{verbatim}
\end{quote}

To use Times Roman, put the following in the preamble:
\begin{quote}
\begin{verbatim}
\usepackage{times}
\end{verbatim}
\end{quote}
(Alternatives like txfonts or newtx are also acceptable.)

Please see the \LaTeX{} source of this document for comments on other packages that may be useful.

Set the title and author using \verb|\title| and \verb|\author|. Within the author list, format multiple authors using \verb|\and| and \verb|\And| and \verb|\AND|; please see the \LaTeX{} source for examples.

By default, the box containing the title and author names is set to the minimum of 5 cm. If you need more space, include the following in the preamble:
\begin{quote}
\begin{verbatim}
\setlength\titlebox{<dim>}
\end{verbatim}
\end{quote}
where \verb|<dim>| is replaced with a length. Do not set this length smaller than 5 cm.

\section{Document Body}

\subsection{Footnotes}

Footnotes are inserted with the \verb|\footnote| command.\footnote{This is a footnote.}

\subsection{Tables and figures}

See Table~\ref{tab:accents} for an example of a table and its caption.
\textbf{Do not override the default caption sizes.}

\begin{table}
  \centering
  \begin{tabular}{lc}
    \hline
    \textbf{Command} & \textbf{Output} \\
    \hline
    \verb|{\"a}|     & {\"a}           \\
    \verb|{\^e}|     & {\^e}           \\
    \verb|{\`i}|     & {\`i}           \\
    \verb|{\.I}|     & {\.I}           \\
    \verb|{\o}|      & {\o}            \\
    \verb|{\'u}|     & {\'u}           \\
    \verb|{\aa}|     & {\aa}           \\\hline
  \end{tabular}
  \begin{tabular}{lc}
    \hline
    \textbf{Command} & \textbf{Output} \\
    \hline
    \verb|{\c c}|    & {\c c}          \\
    \verb|{\u g}|    & {\u g}          \\
    \verb|{\l}|      & {\l}            \\
    \verb|{\~n}|     & {\~n}           \\
    \verb|{\H o}|    & {\H o}          \\
    \verb|{\v r}|    & {\v r}          \\
    \verb|{\ss}|     & {\ss}           \\
    \hline
  \end{tabular}
  \caption{Example commands for accented characters, to be used in, \emph{e.g.}, Bib\TeX{} entries.}
  \label{tab:accents}
\end{table}

As much as possible, fonts in figures should conform
to the document fonts. See Figure~\ref{fig:experiments} for an example of a figure and its caption.

Using the \verb|graphicx| package graphics files can be included within figure
environment at an appropriate point within the text.
The \verb|graphicx| package supports various optional arguments to control the
appearance of the figure.
You must include it explicitly in the \LaTeX{} preamble (after the
\verb|\documentclass| declaration and before \verb|\begin{document}|) using
\verb|\usepackage{graphicx}|.

\begin{figure}[t]
  \includegraphics[width=\columnwidth]{example-image-golden}
  \caption{A figure with a caption that runs for more than one line.
    Example image is usually available through the \texttt{mwe} package
    without even mentioning it in the preamble.}
  \label{fig:experiments}
\end{figure}

\begin{figure*}[t]
  \includegraphics[width=0.48\linewidth]{example-image-a} \hfill
  \includegraphics[width=0.48\linewidth]{example-image-b}
  \caption {A minimal working example to demonstrate how to place
    two images side-by-side.}
\end{figure*}

\subsection{Hyperlinks}

Users of older versions of \LaTeX{} may encounter the following error during compilation:
\begin{quote}
\verb|\pdfendlink| ended up in different nesting level than \verb|\pdfstartlink|.
\end{quote}
This happens when pdf\LaTeX{} is used and a citation splits across a page boundary. The best way to fix this is to upgrade \LaTeX{} to 2018-12-01 or later.

\subsection{Citations}

\begin{table*}
  \centering
  \begin{tabular}{lll}
    \hline
    \textbf{Output}           & \textbf{natbib command} & \textbf{ACL only command} \\
    \hline
    \citep{Gusfield:97}       & \verb|\citep|           &                           \\
    \citealp{Gusfield:97}     & \verb|\citealp|         &                           \\
    \citet{Gusfield:97}       & \verb|\citet|           &                           \\
    \citeyearpar{Gusfield:97} & \verb|\citeyearpar|     &                           \\
    \citeposs{Gusfield:97}    &                         & \verb|\citeposs|          \\
    \hline
  \end{tabular}
  \caption{\label{citation-guide}
    Citation commands supported by the style file.
    The style is based on the natbib package and supports all natbib citation commands.
    It also supports commands defined in previous ACL style files for compatibility.
  }
\end{table*}

Table~\ref{citation-guide} shows the syntax supported by the style files.
We encourage you to use the natbib styles.
You can use the command \verb|\citet| (cite in text) to get ``author (year)'' citations, like this citation to a paper by \citet{Gusfield:97}.
You can use the command \verb|\citep| (cite in parentheses) to get ``(author, year)'' citations \citep{Gusfield:97}.
You can use the command \verb|\citealp| (alternative cite without parentheses) to get ``author, year'' citations, which is useful for using citations within parentheses (e.g. \citealp{Gusfield:97}).

A possessive citation can be made with the command \verb|\citeposs|.
This is not a standard natbib command, so it is generally not compatible
with other style files.

\subsection{References}

\nocite{Ando2005,andrew2007scalable,rasooli-tetrault-2015}

The \LaTeX{} and Bib\TeX{} style files provided roughly follow the American Psychological Association format.
If your own bib file is named \texttt{custom.bib}, then placing the following before any appendices in your \LaTeX{} file will generate the references section for you:
\begin{quote}
\begin{verbatim}
\bibliography{custom}
\end{verbatim}
\end{quote}

You can obtain the complete ACL Anthology as a Bib\TeX{} file from \url{https://aclweb.org/anthology/anthology.bib.gz}.
To include both the Anthology and your own .bib file, use the following instead of the above.
\begin{quote}
\begin{verbatim}
\bibliography{anthology,custom}
\end{verbatim}
\end{quote}

Please see Section~\ref{sec:bibtex} for information on preparing Bib\TeX{} files.

\subsection{Equations}

An example equation is shown below:
\begin{equation}
  \label{eq:example}
  A = \pi r^2
\end{equation}

Labels for equation numbers, sections, subsections, figures and tables
are all defined with the \verb|\label{label}| command and cross references
to them are made with the \verb|\ref{label}| command.

This an example cross-reference to Equation~\ref{eq:example}.

\subsection{Appendices}

Use \verb|\appendix| before any appendix section to switch the section numbering over to letters. See Appendix~\ref{sec:appendix} for an example.

\section{Bib\TeX{} Files}
\label{sec:bibtex}

Unicode cannot be used in Bib\TeX{} entries, and some ways of typing special characters can disrupt Bib\TeX's alphabetization. The recommended way of typing special characters is shown in Table~\ref{tab:accents}.

Please ensure that Bib\TeX{} records contain DOIs or URLs when possible, and for all the ACL materials that you reference.
Use the \verb|doi| field for DOIs and the \verb|url| field for URLs.
If a Bib\TeX{} entry has a URL or DOI field, the paper title in the references section will appear as a hyperlink to the paper, using the hyperref \LaTeX{} package.

\section*{Limitations}

This document does not cover the content requirements for ACL or any
other specific venue.  Check the author instructions for
information on
maximum page lengths, the required ``Limitations'' section,
and so on.

\section*{Acknowledgments}

This document has been adapted
by Steven Bethard, Ryan Cotterell and Rui Yan
from the instructions for earlier ACL and NAACL proceedings, including those for
ACL 2019 by Douwe Kiela and Ivan Vuli\'{c},
NAACL 2019 by Stephanie Lukin and Alla Roskovskaya,
ACL 2018 by Shay Cohen, Kevin Gimpel, and Wei Lu,
NAACL 2018 by Margaret Mitchell and Stephanie Lukin,
Bib\TeX{} suggestions for (NA)ACL 2017/2018 from Jason Eisner,
ACL 2017 by Dan Gildea and Min-Yen Kan,
NAACL 2017 by Margaret Mitchell,
ACL 2012 by Maggie Li and Michael White,
ACL 2010 by Jing-Shin Chang and Philipp Koehn,
ACL 2008 by Johanna D. Moore, Simone Teufel, James Allan, and Sadaoki Furui,
ACL 2005 by Hwee Tou Ng and Kemal Oflazer,
ACL 2002 by Eugene Charniak and Dekang Lin,
and earlier ACL and EACL formats written by several people, including
John Chen, Henry S. Thompson and Donald Walker.
Additional elements were taken from the formatting instructions of the \emph{International Joint Conference on Artificial Intelligence} and the \emph{Conference on Computer Vision and Pattern Recognition}.

% Bibliography entries for the entire Anthology, followed by custom entries
%\bibliography{custom,anthology-overleaf-1,anthology-overleaf-2}

% Custom bibliography entries only
\bibliography{custom}

\appendix

\section{Example Appendix}
\label{sec:appendix}

This is an appendix.

\end{document}
